\section{Introduction}

Reading text in natural images remains a central failure mode for otherwise capable vision-language systems. The TextVQA benchmark was introduced precisely to stress this gap: a successful model must not only parse scene text, but also integrate it with visual context and question semantics to produce a short, exact answer \citep{singh2019textvqa}. In practice, this setting is harder than generic image question answering \citep{antol2015vqa} because many correct responses depend on OCR fidelity, prompt format, answer normalization, and the model's ability to reason over text that is sparse, noisy, or visually embedded in cluttered scenes.

At the same time, the open-source VLM ecosystem has matured enough that a serious empirical study is now feasible without closed commercial models. Architectures such as BLIP-2 \citep{li2023blip2}, LLaVA \citep{liu2023visual}, InternVL2.5 \citep{chen2024internvl25}, and Qwen2.5-VL \citep{qwen2025qwen25vl} offer strong multimodal priors at parameter scales that can still be evaluated and adapted in an academic setting. Yet the existence of capable models does not, by itself, solve the experimental-design problem. A naïve study can easily become either too shallow, by comparing only a small number of zero-shot settings, or too diffuse, by mixing screening, validation, prompt exploration, and fine-tuning without a clear promotion policy. In TextVQA, where OCR choices matter almost as much as backbone choice, such looseness makes the final conclusions difficult to defend.

This paper therefore adopts a deliberately staged protocol rather than a flat benchmark. We first run broad zero-shot screening on a stratified internal development split, then reserve the official validation set for a narrower set of finalists, and finally conduct a winner-conditioned LoRA study on the selected backbone. The key principle is that each stage should answer a different scientific question: screening compares backbones and OCR-sensitive prompting cheaply, finalist reruns verify which combinations are still competitive on the official split, and the training stage isolates the effect of parameter-efficient adaptation while keeping follow-up ablations attached to a selected LoRA family rather than to an arbitrary default. The final claims are made on official validation evidence, while internal-dev results are used for screening, diagnostics, and budget allocation. The result is an experiment that is broader than a single-tuned-model class project, but still structured enough to remain methodologically coherent.

Our study is also designed around realistic compute constraints rather than idealized cluster access. All canonical evaluation runs are executed locally in a resumable queue on a single Apple-Silicon workstation, while only the final LoRA matrix is moved to rented remote GPUs under a separate run namespace. This separation is not merely operational. It preserves a clean experimental boundary between the local zero-shot/finalist evidence and the remote training evidence, reducing the risk that nominally identical runs mix hardware regimes, caches, or checkpoints in ways that complicate interpretation.

The paper makes the following contributions. First, it introduces a reproducible funnel protocol for benchmarking open-source VLMs on TextVQA while preserving the official validation split for finalist and final tuned-model evidence. Second, it treats OCR as a first-class experimental variable across zero-shot prompting, training-time ablation, and post-hoc robustness analysis. Third, it develops a winner-conditioned LoRA study in which training-stage follow-up allocation, internal-dev adapter promotion, and official-validation claims are explicitly separated. This separation keeps a local diagnostic objective distinct from the final paper-facing evaluation objective. Finally, the project delivers a complete implementation pipeline that couples layered configuration, resumable execution, remote multi-GPU orchestration, progress reporting, and generated audit artifacts, making the empirical claims auditable rather than anecdotal.

Section~\ref{sec:related} situates the study with respect to TextVQA, open-source VLMs, and parameter-efficient adaptation. Section~\ref{sec:design} then formalizes the staged experimental protocol and systems boundary that govern the full study.
